\documentclass[11pt]{article}
\usepackage[a4paper,margin=2.5cm]{geometry}
\usepackage[T1]{fontenc}\usepackage{lmodern}\usepackage[hidelinks]{hyperref}
\setlength{\parindent}{0pt}\setlength{\parskip}{6pt}
\begin{document}
Cong-Danh Tran\\
University of Information Technology, Ho Chi Minh City, Viet Nam\\
Vietnam National University, Ho Chi Minh City, Viet Nam\\
danhtc@grad.uit.edu.vn \quad ORCID 0009-0001-9660-8808

\today

The Editor-in-Chief\\
International Journal of Medical Informatics

Dear Editor,

I submit the original research article \emph{``Attribution is not mitigation: a paired-label counterfactual audit for hypertension clinical decision support''} for consideration in the International Journal of Medical Informatics.

Counterfactual explanations are increasingly proposed as patient- or clinician-facing outputs of decision support models. This study asks whether the way a training label is recorded changes these explanations, and whether the common remedy of locking implausible features removes the effect. Using two hypertension labels recorded for the same NHANES participants (self-reported diagnosis and examination-based blood pressure), the study shows that diagnosis-related features contribute to the difference between the two models' responses, but that locking these features does not reduce it consistently, and that guideline-based constraints leave valid counterfactuals for few people. The proposed Paired-Label Counterfactual Audit (PLCA) is intended as a pre-deployment audit step for developers and governance teams; the manuscript includes a short clinical translation roadmap and states plainly that clinical impact has not been evaluated.

The work builds on a recent article in this journal (Thieu, Int.\ J.\ Med.\ Inform.\ 219 (2026) 106555), which suggested that self-reported labels may cause implausible counterfactuals but had no measured comparison label; PLCA tests this question with two labels for the same people.

The main text has \WORDS{} words, with 4 figures and 5 tables; model details, performance and calibration, and data characteristics are in Appendix A. The completed IJMEDI checklist and the TRIPOD+AI checklist are provided. The data are public, and all code and results are available at \url{https://github.com/danh-tc/plca-hypertension}. Generative AI use is declared in the manuscript.

This manuscript has not been published and is not under consideration elsewhere. I am the sole author and declare no competing interests.

Sincerely,\\[4pt]
Cong-Danh Tran
\end{document}
